\clearpage\thispagestyle{empty}
\pdfbookmark[0]{Quellen zur Besprechung von d’Aguila}{aguila1812-sources}
\begin{center}\textbf{Quellen zur Besprechung von d’Aguila}\end{center}
\noindent Diese KI-Übersetzung ins heutige Deutsch folgt Carl Friedrich Gauß, \textit{Werke}, Bd.~VI (Göttingen, 1874), S.~541. Die Besprechung erschien zuerst unter Paris in den \textit{Göttingischen gelehrten Anzeigen}, Stück~102, 27. Juni 1812, S.~1009--1010.
\par\smallskip\noindent Die erste Ziffer des Nenners ist im Druck der Werkausgabe schwer zu erkennen. Der Erstdruck auf S.~1009 stützt die Lesung \(1/59\). Die Bibliothèque nationale de France verzeichnet das besprochene \href{https://catalogue.bnf.fr/ark:/12148/cb30005371x}{Buch von C. J. E. H. d’Aguila} als 1806 in Paris bei Delance gedruckt. In der Werkausgabe steht der Name des Druckers als \textit{Délance}.
\par\smallskip\noindent \href{https://gdz.sub.uni-goettingen.de/id/PPN236007467}{Text der Werkausgabe}. Erstdruck, \href{https://images.sub.uni-goettingen.de/iiif/image/gdz:PPN319721507_1812:00001049/full/full/0/default.jpg}{S.~1009} und \href{https://images.sub.uni-goettingen.de/iiif/image/gdz:PPN319721507_1812:00001050/full/full/0/default.jpg}{S.~1010}.
\par\smallskip\noindent Übersetzung und Satz wurden von KI erstellt.
\clearpage
